\documentclass[preview]{standalone}

\usepackage{amsmath}
\usepackage{amssymb}
\usepackage{stellar}
\usepackage{definitions}
\usepackage{bettelini}

\begin{document}

\id{probability-conditional-distributions}
\genpage

\section{Discrete conditional laws}

\begin{snippetdefinition}{conditional-law-discrete-definition}{Conditional law in the discrete case}
    Let \((\Omega,\powerset(\Omega),\mathbb P)\) be a
    \countableprobabilityspace. Let
    \(X\colon\Omega\fromto E\) and \(Y\colon\Omega\fromto F\) be
    \discreterandomvariable[discrete random variables], where
    \(E,F\subseteq\realnumbers\) are finite or countably infinite
    \set[sets]. If \(y\in F\) satisfies \(\mathbb P(Y=y)>0\), the
    \emph{conditional law} of \(X\) given \(Y=y\) is the
    \probabilitymeasure on \((E,\powerset(E))\) defined by
    \[
        \mu_{X\mid Y=y}(A)
        =
        \mathbb P(X\in A\mid Y=y)
        =
        \frac{\mathbb P(X\in A,\ Y=y)}{\mathbb P(Y=y)}
    \]
    for every \(A\subseteq E\).
\end{snippetdefinition}

\begin{snippetdefinition}{conditional-discrete-density-definition}{Conditional discrete density}
    Let \((\Omega,\powerset(\Omega),\mathbb P)\) be a
    \countableprobabilityspace. Let
    \(X\colon\Omega\fromto E\) and \(Y\colon\Omega\fromto F\) be
    \discreterandomvariable[discrete random variables], where
    \(E,F\subseteq\realnumbers\) are finite or countably infinite
    \set[sets]. Let \(p_{X,Y}\) be the \jointdensity of \((X,Y)\), and let
    \(p_Y\) be the \randomvariabledensity of \(Y\). If \(y\in F\) satisfies
    \(p_Y(y)>0\), the \emph{conditional discrete density} of \(X\) given
    \(Y=y\) is
    \[
        p_{X\mid Y}(x\mid y)
        =
        \frac{p_{X,Y}(x,y)}{p_Y(y)},
        \qquad x\in E.
    \]
\end{snippetdefinition}

\begin{snippetproposition}{conditional-discrete-density-is-density}{Conditional discrete densities are densities}
    Let \((\Omega,\powerset(\Omega),\mathbb P)\) be a
    \countableprobabilityspace. Let
    \(X\colon\Omega\fromto E\) and \(Y\colon\Omega\fromto F\) be
    \discreterandomvariable[discrete random variables], where
    \(E,F\subseteq\realnumbers\) are finite or countably infinite
    \set[sets]. If \(y\in F\) satisfies \(p_Y(y)>0\), then
    \(p_{X\mid Y}(\cdot\mid y)\) is a \discretedensity on \(E\), and
    \[
        \mu_{X\mid Y=y}(A)
        =
        \sum_{x\in A}p_{X\mid Y}(x\mid y)
    \]
    for every \(A\subseteq E\).
\end{snippetproposition}

\begin{snippetproof}{conditional-discrete-density-is-density-proof}{conditional-discrete-density-is-density}{Conditional discrete densities are densities}
    Nonnegativity is immediate. Moreover,
    \[
        \sum_{x\in E}p_{X\mid Y}(x\mid y)
        =
        \frac{1}{p_Y(y)}
        \sum_{x\in E}p_{X,Y}(x,y)
        =
        \frac{p_Y(y)}{p_Y(y)}
        =
        1.
    \]
    For \(A\subseteq E\),
    \[
        \sum_{x\in A}p_{X\mid Y}(x\mid y)
        =
        \frac{1}{p_Y(y)}
        \sum_{x\in A}\mathbb P(X=x,Y=y)
        =
        \frac{\mathbb P(X\in A,Y=y)}{\mathbb P(Y=y)}.
    \]
\end{snippetproof}

\begin{snippetproposition}{conditional-discrete-density-independent-case}{Conditional density under independence}
    Let \((\Omega,\powerset(\Omega),\mathbb P)\) be a
    \countableprobabilityspace. Let
    \(X\colon\Omega\fromto E\) and \(Y\colon\Omega\fromto F\) be independent
    \discreterandomvariable[discrete random variables], where
    \(E,F\subseteq\realnumbers\) are finite or countably infinite
    \set[sets]. If \(y\in F\) satisfies \(p_Y(y)>0\), then
    \[
        p_{X\mid Y}(x\mid y)=p_X(x)
    \]
    for every \(x\in E\).
\end{snippetproposition}

\begin{snippetproof}{conditional-discrete-density-independent-case-proof}{conditional-discrete-density-independent-case}{Conditional density under independence}
    By independence,
    \[
        p_{X,Y}(x,y)=p_X(x)p_Y(y).
    \]
    Since \(p_Y(y)>0\),
    \[
        p_{X\mid Y}(x\mid y)
        =
        \frac{p_X(x)p_Y(y)}{p_Y(y)}
        =
        p_X(x).
    \]
\end{snippetproof}

\begin{snippetdefinition}{discrete-conditional-expectation-definition}{Conditional expected value in the discrete case}
    Let \((\Omega,\powerset(\Omega),\mathbb P)\) be a
    \countableprobabilityspace. Let
    \(X\colon\Omega\fromto E\subseteq\realnumbers\) and
    \(Y\colon\Omega\fromto F\subseteq\realnumbers\) be
    \discreterandomvariable[discrete random variables], where \(E\) and
    \(F\) are finite or countably infinite \set[sets]. If \(y\in F\)
    satisfies \(p_Y(y)>0\) and
    \[
        \sum_{x\in E}|x|p_{X\mid Y}(x\mid y)<+\infty,
    \]
    then the \emph{conditional expected value} of \(X\) given \(Y=y\) is
    \[
        \expectation[X\mid Y=y]
        =
        \sum_{x\in E}xp_{X\mid Y}(x\mid y).
    \]
    If the above sum is finite for every \(y\in F\) with \(p_Y(y)>0\), the
    elementary conditional expectation of \(X\) given \(Y\) is the
    \randomvariable
    \[
        \expectation[X\mid Y]
        =
        g(Y),
        \qquad
        g(y)=\expectation[X\mid Y=y].
    \]
\end{snippetdefinition}

\begin{snippetexample}{poisson-split-conditional-binomial-example}{Conditioning two independent Poisson variables on their sum}
    Let \(X\) and \(Y\) be independent real-valued
    \discreterandomvariable[discrete random variables] with
    \[
        X\sim\operatorname{Poisson}(\lambda),
        \qquad
        Y\sim\operatorname{Poisson}(\gamma),
    \]
    where \(\lambda,\gamma>0\). Let \(n\in\naturalnumbers\). Then, for
    \(k\in\{0,\ldots,n\}\),
    \[
        \mathbb P(X=k\mid X+Y=n)
        =
        \binom{n}{k}
        \left(\frac{\lambda}{\lambda+\gamma}\right)^k
        \left(\frac{\gamma}{\lambda+\gamma}\right)^{n-k}.
    \]
    Therefore
    \[
        X\mid X+Y=n
        \sim
        \operatorname{Binomial}\left(n,\frac{\lambda}{\lambda+\gamma}\right),
    \]
    and
    \[
        \expectation[X\mid X+Y=n]
        =
        n\frac{\lambda}{\lambda+\gamma}.
    \]
\end{snippetexample}

\begin{snippetproof}{poisson-split-conditional-binomial-example-proof}{poisson-split-conditional-binomial-example}{Conditioning two independent Poisson variables on their sum}
    Since \(X+Y\sim\operatorname{Poisson}(\lambda+\gamma)\), for
    \(k\in\{0,\ldots,n\}\),
    \begin{align*}
        \mathbb P(X=k\mid X+Y=n)
        &=
        \frac{\mathbb P(X=k,\ X+Y=n)}
        {\mathbb P(X+Y=n)}
        \\
        &=
        \frac{\mathbb P(X=k,\ Y=n-k)}
        {\mathbb P(X+Y=n)}
        \\
        &=
        \frac{
            e^{-\lambda}\frac{\lambda^k}{k!}
            e^{-\gamma}\frac{\gamma^{n-k}}{(n-k)!}
        }{
            e^{-(\lambda+\gamma)}
            \frac{(\lambda+\gamma)^n}{n!}
        }
        \\
        &=
        \binom{n}{k}
        \left(\frac{\lambda}{\lambda+\gamma}\right)^k
        \left(\frac{\gamma}{\lambda+\gamma}\right)^{n-k}.
    \end{align*}
    This is the density of the stated \binomialdistribution. Its expected
    value is the mean of a binomial law.
\end{snippetproof}

\section{Continuous conditional densities}

\begin{snippetnote}{continuous-conditioning-on-null-event-note}{Conditioning on a point in the continuous case}
    Let \((X,Y)\) be an \absolutelycontinuousrandomvector in
    \(\realnumbers^2\). Then typically
    \[
        \mathbb P(Y=y)=0
    \]
    for every fixed \(y\in\realnumbers\). Thus the expression
    \(\mathbb P(X\in A\mid Y=y)\) is not ordinary
    \conditionalprobability given an \event of positive probability. In the
    density setting, conditioning on \(Y=y\) is represented by a regular
    conditional density.
\end{snippetnote}

\begin{snippetdefinition}{conditional-density-absolutely-continuous-definition}{Conditional density in the absolutely continuous case}
    Let \((\Omega,\mathcal F,\mathbb P)\) be a \probabilityspace, and let
    \((X,Y)\) be an \absolutelycontinuousrandomvector in \(\realnumbers^2\)
    with \continuousjointdensity \(f_{X,Y}\). Let
    \[
        f_Y(y)=\int_\realnumbers f_{X,Y}(x,y)\,dx
    \]
    be the \continuousmarginaldensity of \(Y\). For every
    \(y\in\realnumbers\) such that \(f_Y(y)>0\), the
    \emph{conditional density} of \(X\) given \(Y=y\) is
    \[
        f_{X\mid Y=y}(x)
        =
        f_{X\mid Y}(x\mid y)
        =
        \frac{f_{X,Y}(x,y)}{f_Y(y)}
    \]
    for almost every \(x\in\realnumbers\).
\end{snippetdefinition}

\begin{snippetproposition}{conditional-continuous-density-is-density}{Conditional continuous densities are densities}
    Let \((\Omega,\mathcal F,\mathbb P)\) be a \probabilityspace, and let
    \((X,Y)\) be an \absolutelycontinuousrandomvector in \(\realnumbers^2\)
    with \continuousjointdensity \(f_{X,Y}\). Let
    \[
        f_Y(y)=\int_\realnumbers f_{X,Y}(x,y)\,dx
    \]
    be the \continuousmarginaldensity of \(Y\). If \(y\in\realnumbers\)
    satisfies \(f_Y(y)>0\), then
    \[
        f_{X\mid Y=y}(x)
        =
        \frac{f_{X,Y}(x,y)}{f_Y(y)}
    \]
    is a \probabilitydensity on \(\realnumbers\).
\end{snippetproposition}

\begin{snippetproof}{conditional-continuous-density-is-density-proof}{conditional-continuous-density-is-density}{Conditional continuous densities are densities}
    Nonnegativity follows from nonnegativity of \(f_{X,Y}\). Moreover,
    \[
        \int_\realnumbers f_{X\mid Y=y}(x)\,dx
        =
        \frac{1}{f_Y(y)}
        \int_\realnumbers f_{X,Y}(x,y)\,dx
        =
        1.
    \]
\end{snippetproof}

\begin{snippetproposition}{conditional-continuous-density-independent-case}{Conditional continuous density under independence}
    Let \((\Omega,\mathcal F,\mathbb P)\) be a \probabilityspace, and let
    \((X,Y)\) be an \absolutelycontinuousrandomvector in \(\realnumbers^2\)
    with \continuousjointdensity \(f_{X,Y}\). Let \(f_X\) and \(f_Y\) be the
    \continuousmarginaldensity[marginal densities] obtained by integrating
    \(f_{X,Y}\). If \(X\) and \(Y\) are independent, then, for every \(y\)
    with \(f_Y(y)>0\),
    \[
        f_{X\mid Y}(x\mid y)=f_X(x)
    \]
    for almost every \(x\in\realnumbers\).
\end{snippetproposition}

\begin{snippetproof}{conditional-continuous-density-independent-case-proof}{conditional-continuous-density-independent-case}{Conditional continuous density under independence}
    By the density factorization criterion,
    \[
        f_{X,Y}(x,y)=f_X(x)f_Y(y)
    \]
    almost everywhere. Dividing by \(f_Y(y)>0\) gives the formula.
\end{snippetproof}

\begin{snippetdefinition}{continuous-conditional-expectation-definition}{Conditional expected value through a conditional density}
    Let \((\Omega,\mathcal F,\mathbb P)\) be a \probabilityspace, and let
    \((X,Y)\) be an \absolutelycontinuousrandomvector in \(\realnumbers^2\)
    with \conditionaldensity \(f_{X\mid Y=y}\). If
    \[
        \int_\realnumbers |x|f_{X\mid Y=y}(x)\,dx<+\infty,
    \]
    then
    \[
        \expectation[X\mid Y=y]
        =
        \int_\realnumbers x f_{X\mid Y=y}(x)\,dx.
    \]
\end{snippetdefinition}

\begin{snippetexample}{continuous-conditional-exponential-example}{A continuous conditional exponential law}
    Let \((X,Y)\) be an \absolutelycontinuousrandomvector in
    \(\realnumbers^2\) with density
    \[
        f_{X,Y}(x,y)
        =
        \frac{e^{-x/y}e^{-y}}{y}\indicator_{(0,+\infty)^2}(x,y).
    \]
    Then \(Y\sim\operatorname{Exp}(1)\), and for every \(y>0\),
    \[
        X\mid Y=y\sim\operatorname{Exp}\left(\frac{1}{y}\right).
    \]
    In particular,
    \[
        \expectation[X\mid Y=y]=y.
    \]
\end{snippetexample}

\begin{snippetproof}{continuous-conditional-exponential-example-proof}{continuous-conditional-exponential-example}{A continuous conditional exponential law}
    For \(y>0\),
    \begin{align*}
        f_Y(y)
        &=
        \int_\realnumbers
        \frac{e^{-x/y}e^{-y}}{y}\indicator_{(0,+\infty)^2}(x,y)\,dx
        \\
        &=
        e^{-y}\int_0^{+\infty}\frac{e^{-x/y}}{y}\,dx
        =
        e^{-y}.
    \end{align*}
    For \(y\leq0\), \(f_Y(y)=0\). Thus
    \(Y\sim\operatorname{Exp}(1)\). For \(y>0\),
    \[
        f_{X\mid Y=y}(x)
        =
        \frac{f_{X,Y}(x,y)}{f_Y(y)}
        =
        \frac{1}{y}e^{-x/y}\indicator_{(0,+\infty)}(x).
    \]
    This is the density of
    \(\operatorname{Exp}\left(\frac{1}{y}\right)\), whose mean is \(y\).
\end{snippetproof}

\end{document}
